\section{Introduction}

Crops grown aboard the International Space Station (ISS) in the Veggie plant
growth system are intended to supplement crew diets on long-duration missions,
supplying fresh produce that resupply cannot economically provide
\citep{massa2017}. Whether they do so usefully depends on what is actually in
the tissue. Spaceflight alters plant water relations, root-zone oxygen
availability and nutrient transport, and there is a reasonable prior that
mineral accumulation differs from ground controls.

The same minerals matter on the crew side. Spaceflight perturbs human mineral
homeostasis: bone resorption releases calcium, fluid shifts alter electrolyte
balance, and iron stores rise as erythropoiesis is suppressed
\citep{smith2005,zheng2023}. If spaceflight-grown produce were to concentrate
particular ions, a plant-rich diet could in principle interact with those
perturbations --- an accumulation-plus-overload concern that motivates asking
what the plants actually contain and what a realistic serving delivers.

Answering that requires two things this literature has lacked. The first is
harmonised elemental data across more than one crop. Two NASA Open Science Data
Repository (OSDR) studies report ICP-OES elemental composition for Veggie
produce: OSD-745 for red romaine lettuce \citep{khodadad2020,osdr745} and
OSD-655 for Tokyo Bekana \citep{bunchek2024,osdr655}. They have not previously
been analysed together, and as this work shows, doing so naively produces a
result that is wrong by three orders of magnitude for five elements.

The second is a clear picture of what human spaceflight mineral and dietary data
exists and can be obtained. The NASA Life Sciences Data Archive (LSDA) holds
records for several thousand experiments, but its catalog is designed for
interactive browsing rather than systematic survey, and an experiment's record
does not make its accessibility obvious at a glance.

\begin{smamethod}[title=Contributions]
\begin{enumerate}\itemsep2pt
  \item A machine-readable survey of all \LSDATotal{} LSDA experiments,
        annotated for mineral and nutritional content and classified by whether
        the data can actually be obtained.
  \item Identification and correction of a unit-labelling defect in OSD-655
        that inflates five elements by a factor of 1000, present in the archived
        data dictionary as well as the data files.
  \item A compositional (centred log-ratio) reanalysis separating genuine
        ionome changes from dry-matter dilution artefacts, showing that the two
        Veggie crops respond differently to spaceflight.
  \item A quantitative assessment of what a Veggie serving contributes against
        NASA crew mineral requirements \citep{jsc67378}, and of whether dietary
        mineral overload is achievable.
\end{enumerate}
\end{smamethod}

We state at the outset what this work does not do. The plant and crew datasets
come from different missions and different people; nothing here supports a
causal claim that Veggie produce altered anyone's serum chemistry. That question
requires paired dietary-intake and serum data, which exists --- as the
Nutritional Status Assessment (SMO-016E) --- but is not openly available.
Section~\ref{sec:availability} sets out the route to it.
